Document compliance is a crucial process in company environments, certifying that every form of document follows regulatory standards. This process ensures that every document within this scenario follows a series of human and automatic verifications, and one of these steps is document classification. In most cases, document classification follows a simple and traditional classification: a model is trained with examples of every class, and later, categorizes documents into these classes.
The challenge arises when these documents update their structure. In this case, the traditional classification model must be retrained, as the new document type is not mapped in the current model. Therefore, a possible solution to this problem is to train the model under the \gls{ZSL} paradigm. This means that the model will be able to classify completely unseen documents, making it more resilient to changes. Therefore, this work tackles the \gls{DIC} problem in a \gls{ZSL} paradigm, introducing the \gls{VDM} approach, which frames the classification task as a matching problem. To achieve such a model, this work also presents the \gls{LA-CDIP} dataset, a dataset specialized in \gls{ZSL}-\gls{DIC}, providing ample document layout diversity to allow great generalization in models trained from scratch. This dataset has 17,295 documents, with 665 different classes, where every class is a different layout arrangement.
This dataset is constructed with a \gls{DL} assisted framework, where a \gls{DL} model is trained with an intermediary dataset, and in return, uses its acquired knowledge to accelerate the human labeling process.
To further stimulate research with this framework, this work also provides an extensive benchmark, taking into consideration multiple \gls{DL} backbones and hyperparameters. In zero-shot scenarios, the proposed method with the new dataset achieves an \acrfull{EER} below 5\% in 1-vs-1 verification with cross-validation, proving to be an efficient solution to the \gls{ZSL}-\gls{DIC} problem.

% O \emph{abstract} é o resumo feito na língua Inglesa. Embora o conteúdo apresentado
% deva ser o mesmo, este texto não deve ser a tradução literal de cada palavra ou
% frase do resumo, muito menos feito em um tradutor automático. É uma língua
% diferente e o texto deveria ser escrito de acordo com suas nuances (aproveite para ler
% \url{http://dx.doi.org/10.6061%2Fclinics%2F2014(03)01}). Por exemplo: \emph{This work presents useful information on how to create a scientific text to describe
% and provide examples of how to use the Computer Science Department's \LaTeX\ class. The \unbcic\
% class defines a standard format for texts, simplifying the process of generating
% CIC documents and enabling authors to focus only on content. The standard was approved
% by the Department's professors and used to create this document. Future work includes
% continued support for the class and improvements on the explanatory text.}